% subsection: 総和型

\subsection{総和型}\label{節：総和型}
  \begin{bookpoint}
    \textbf{和を消す}\par
    添字を一つずらして引けば,総和の増分が一項として現れる.引き算で落とした\bookterm{initial}{初期条件}も忘れずに戻す.
  \end{bookpoint}
  \subsubsection{\texorpdfstring{和と項の関係 \, $S_n=pa_n+f(n)$}{和と項の関係}}
    \bookterm{sequence}{数列}$a_n$の和$S_n=\sum_{k=1}^n a_k$が,
    \[
      S_n=pa_n+f(n)
    \]
    を満たす場合を考えよう.$S_{n+1}-S_n=a_{n+1}$を使えば,$a_n$だけの\bookterm{recurrence}{漸化式}にできる.
    ただし,差を取ると定数分の情報が消えるので,先に$n=1$の条件
    \[
      a_1=pa_1+f(1),\qquad (1-p)a_1=f(1)
    \]
    を残しておく.$p\ne1$のとき,
    \begin{align*}
      S_n&=pa_n+f(n) \booktag{sum-relation-1}{1}\\
      S_{n+1}&=pa_{n+1}+f(n+1) \booktag{sum-relation-2}{2}
    \end{align*}
    の$\bookeqbackref{sum-relation-2}{(2)}-\bookeqbackref{sum-relation-1}{(1)}$から,
    \begin{align*}
      a_{n+1}&=pa_{n+1}-pa_n+f(n+1)-f(n),\\
      a_{n+1}&=\frac p{p-1}a_n+\frac{f(n+1)-f(n)}{1-p}
    \end{align*}
    を得る.付加項を$g(n)=\dfrac{f(n+1)-f(n)}{1-p}$と書けば,$a_{n+1}=\frac p{p-1}a_n+g(n)$である.
    多項式や指数の付加項なら,\booknavref{節：階差等比複合型}{付加項を引く・指数で割る方法}を使える.\footnote{係数$\dfrac{p}{p-1}$は決して$1$にならない.}
    それ以外も,係数が非零なら\booknavref{節：変数係数の一次漸化式}{倍率をそろえて和を取る方法}で表示できる.$p=0$なら次項を直接求めればよい.

    $p=1$では割れない.$S_n=a_n+f(n)$と$S_n=S_{n-1}+a_n$を比較すると,
    \[
      a_{n-1}=f(n)-f(n-1)\quad(n\ge2)
    \]
    となる.$n-1$を$n$と読み替えれば,
    \[
      a_n=f(n+1)-f(n)\quad(n\ge1)
    \]
    を得る.元の$n=1$の式から,$f(1)=0$も必要である.

    逆に,得られた漸化式と最初の条件から元の和の関係を戻せる.
    $T_n=S_n-pa_n-f(n)$と置くと,差の式は$T_{n+1}-T_n=0$を意味する.
    最初の条件が$T_1=0$を与えるので,すべての$n$で$T_n=0$となる.
    \textbf{差を求めた後には,初めの値が必要になる.}

    例題(3)では二つを選ぶ個数を表す記号
    \[
      \binom n2=\frac{n(n-1)}2
    \]
    を使う.$n=1$では二つを選べないので,$\binom12=0$とする.\footnote{この記号を\booknavref{節：二項定理}{二項係数}という.一般の定義と公式は参照先にある.}
    \noindent \bookblocklink{example-summation-1}{problem}{answer}{【例題】}\par\nobreak
    \begin{enumerate}[label=(\arabic*),parsep=3pt,format=\bookitemlink{example-summation-1}{problem}{answer}]
      \item $\displaystyle S_n=2a_n-n$
      \item $\displaystyle S_n=\frac12a_n+n$
      \item $\displaystyle S_n=a_n+\binom n2$
    \end{enumerate}

    \noindent\bookblocklink{example-summation-1}{hint}{problem}{【方針】}\par\nobreak
    \begin{enumerate}[label=(\arabic*),parsep=3pt,format=\bookitemlink{example-summation-1}{hint}{problem}]
      \item まず$n=1$を代入して$a_1$を求めてから,本文の変形をなぞって$a_{n+1}$と$a_n$だけの\bookterm{recurrence}{漸化式}を作りましょう.
      \item 同様です.できた\bookterm{recurrence}{漸化式}では,不動点を引いて定数項を消します.
      \item $p=1$の場合です.$S_n-S_{n-1}=a_n$と比較して,$a_{n-1}$を直接求めましょう.
    \end{enumerate}

    \noindent\bookblocklink{example-summation-1}{answer}{problem}{【解答】}\par\nobreak
    \begin{enumerate}[label=(\arabic*),parsep=3pt,format=\bookitemlink{example-summation-1}{answer}{problem}]
      \item $\displaystyle a_n=2^n-1$
      \item $\displaystyle a_n=1+(-1)^{n-1}$
      \item $\displaystyle a_n=n$
    \end{enumerate}

    \noindent\bookblocklink{example-summation-1}{solution}{problem}{【解法】}\par\nobreak
    \begin{enumerate}[label=(\arabic*),parsep=3pt,format=\bookitemlink{example-summation-1}{solution}{problem}]
    \item $\displaystyle S_n=2a_n-n$\\
        $n=1$の式$a_1=2a_1-1$から,$a_1=1$である.差を取ると,
        \begin{align*}
          a_{n+1}&=\{2a_{n+1}-(n+1)\}-(2a_n-n)\\
          &=2a_{n+1}-2a_n-1
        \end{align*}
        より,$a_{n+1}=2a_n+1$となる.$b_n=a_n+1$と置けば,
        \[
          b_{n+1}=2b_n,\quad b_1=2
        \]
        なので,$b_n=2^n$,従って$a_n=2^n-1$を得る.
    \item $\displaystyle S_n=\frac12a_n+n$\\
        $n=1$の式$a_1=\frac12a_1+1$から,$a_1=2$である.
        本文の公式に$p=\dfrac{1}{2},\ f(n)=n$を入れると,$a_{n+1}=-a_n+2$となる.
        不動点$1$を引き,$b_n=a_n-1$と置けば,
        \[
          b_{n+1}=-b_n,\quad b_1=1
        \]
        なので,$b_n=(-1)^{n-1}$,従って$a_n=1+(-1)^{n-1}$を得る.
    \item $\displaystyle S_n=a_n+\binom n2$\\
        $n=1$では$a_1=a_1+0$となり,この式だけでは$a_1$は定まらない.
        $n\ge2$で差を取ると,
        \begin{align*}
          a_n&=S_n-S_{n-1}\\
          &=a_n-a_{n-1}+\binom n2-\binom{n-1}2\\
          &=a_n-a_{n-1}+n-1
        \end{align*}
        より,$a_{n-1}=n-1$を得る.$n-1$を$n$と読み替えれば,$a_n=n$は$n\ge1$で成り立つ.
    \end{enumerate}

  \subsubsection{\texorpdfstring{和の漸化式 \, $S_{n+1}=pS_n+f(n) $}{和の漸化式}}
    両辺が総和でも,添字をずらして引けば一項ずつの式になる.
    \begin{align*}
      S_{n+1}&=pS_n+f(n) \booktag{sum-pure-1}{1}\\
      S_n&=pS_{n-1}+f(n-1) \booktag{sum-pure-2}{2}
    \end{align*}
    $\bookeqbackref{sum-pure-1}{(1)}-\bookeqbackref{sum-pure-2}{(2)}$から,
    \[
      a_{n+1}=pa_n+f(n)-f(n-1)\quad(n\ge2)
    \]
    を得る.\bookeqref{sum-pure-2}{(2)}は\bookeqref{sum-pure-1}{(1)}の添字を一つ戻した式なので,$n=1$には使えない.
    $a_1=S_1$を保ち,
    \[
      a_2=S_2-S_1=(p-1)S_1+f(1)
    \]
    を別に求め,$a_2$から解く.
    あとは\bookterm{difference}{階差}等比複合型の計算である($p=1$なら\bookterm{difference}{階差}\bookterm{sequence}{数列}型).
    $S_n$を先に解き,$a_1=S_1,\ a_n=S_n-S_{n-1}$($n\ge2$)と戻してもよい.
    \textbf{$S_n$と$a_n$のどちらが解きやすいかで選ぶ.}

    \noindent \bookblocklink{example-summation-2}{problem}{answer}{【例題】}\par\nobreak
    \begin{enumerate}[label=(\arabic*),parsep=3pt,format=\bookitemlink{example-summation-2}{problem}{answer}]
      \item $\displaystyle S_1=1,\ S_{n+1}=2S_n+n$
      \item $\displaystyle S_1=0,\ S_{n+1}=2S_n+n$
    \end{enumerate}

    \noindent\bookblocklink{example-summation-2}{hint}{problem}{【方針】}\par\nobreak
    \begin{enumerate}[label=(\arabic*),parsep=3pt,format=\bookitemlink{example-summation-2}{hint}{problem}]
      \item 本文の式は$n\ge2$でしか使えないことに注意し,$a_2$だけは$S_2=2S_1+1$から別に求めましょう.
      \item 同様に$n\ge2$で\bookterm{recurrence}{漸化式}を作り,$a_2$から解きます.$a_1$も含む\bookterm{general}{一般項}になるか最後に確認しましょう.
    \end{enumerate}

    \noindent\bookblocklink{example-summation-2}{answer}{problem}{【解答】}\par\nobreak
    \begin{enumerate}[label=(\arabic*),parsep=3pt,format=\bookitemlink{example-summation-2}{answer}{problem}]
      \item $\displaystyle a_1=1,\quad a_n=3\cdot2^{n-2}-1\ (n\ge2)$
      \item $\displaystyle a_n=2^{n-1}-1\ (n\ge1)$
    \end{enumerate}

    \noindent\bookblocklink{example-summation-2}{solution}{problem}{【解法】}\par\nobreak
    \begin{enumerate}[label=(\arabic*),parsep=3pt,format=\bookitemlink{example-summation-2}{solution}{problem}]
    \item $\displaystyle S_1=1,\ S_{n+1}=2S_n+n$\\
        $a_1=S_1=1,\ S_2=2S_1+1=3$から,$a_2=2$である.
        本文の公式に$p=2,\ f(n)=n$を入れると,
        \[
          a_{n+1}=2a_n+1\quad(n\ge2)
        \]
        となる.$b_n=a_n+1$と置けば,$b_{n+1}=2b_n,\ b_2=3$なので,
        \[
          a_1=1,\qquad a_n=3\cdot2^{n-2}-1\quad(n\ge2)
        \]
        を得る.この一般項の式は$n=1$には使えない.
    \item $\displaystyle S_1=0,\ S_{n+1}=2S_n+n$\\
        $a_1=0,\ S_2=1$から,$a_2=1$である.
        (1)と同じ更新$a_{n+1}=2a_n+1$($n\ge2$)を,$b_2=a_2+1=2$から解けば,
        \[
          a_n=2^{n-1}-1\quad(n\ge2)
        \]
        となる.$n=1$でも$2^0-1=0=a_1$なので,今回は同じ式が$n\ge1$で成り立つ.
    \end{enumerate}
